\chapter{Myasthenia Gravis Tests}
\section*{What Is This Test?} Myasthenia gravis testing evaluates fluctuating weakness caused by impaired nerve-to-muscle transmission.\section*{Alternative Names} AChR antibody test; MuSK antibody test; neuromuscular-junction testing.\section*{Principle} Antibody assays and electrodiagnostic tests detect immune or physiologic evidence of impaired transmission.\section*{Why This Test Is Ordered} It is used for variable eyelid droop, double vision, chewing or swallowing fatigue, limb weakness, or respiratory weakness.\section*{Primary vs Secondary Test} Neurologic examination is primary; antibodies, repetitive-nerve stimulation, and single-fiber EMG support diagnosis.\section*{How This Test Is Performed} Blood is tested for AChR, MuSK, or other antibodies; nerve and muscle responses may be measured.\section*{What Preparation Is Needed} Usually none; medication instructions for electrodiagnostic testing are individualized.\section*{Understanding Results} A positive specific antibody supports the diagnosis, but some patients are seronegative. Results must match the clinical pattern.\section*{Follow-up Tests} Chest CT for thymic disease, pulmonary-function testing, additional antibodies, and specialist assessment may follow. New breathing or swallowing difficulty is urgent.\section*{References} \begin{enumerate}\item MedlinePlus. Myasthenia Gravis Tests.\item International Consensus Guidance for Management of Myasthenia Gravis.\end{enumerate}\clearpage\section*{Understanding Results and Follow-up}
The laboratory report should identify the specimen, method, units, reference interval or decision limit, and whether the result is positive, negative, abnormal, or inconclusive. Reference intervals vary with age, sex, pregnancy, specimen type, assay, and laboratory; a result outside a range is not automatically a diagnosis, and a result within a range does not always exclude disease. Discuss unexpected results with the ordering clinician, who can decide whether repeat or confirmatory testing, treatment, or referral is appropriate. Do not change medicines or delay urgent care based on an isolated result.
\section*{Regulatory Status and Authorization Date}This chapter describes a test category, not a specific commercial product. FDA, CDSCO, EU IVDR, and MHRA approval or registration dates therefore cannot be assigned without the assay manufacturer, product name, instrument, and intended use. For laboratory-developed tests, an individual agency approval date may not exist. See the Preface for the interpretation of regulatory status.
\clearpage
